\sectionkicker{Source-backed technical notes}
\subsection*{Frontier Models — モデル競争から利用形態の競争へ: Technical Notes}
本欄は記事本文で圧縮した一次資料上の情報を、比較・再検証しやすい形で整理したものである。時系列、確認済みの主張、留意点、一次資料URLを掲載する。
\medskip
\subsection*{Theme at a glance}
\addcontentsline{toc}{subsection}{Theme at a glance}
\begin{center}
\footnotesize
\begin{tabularx}{\linewidth}{@{}p{0.20\linewidth}p{0.10\linewidth}p{0.14\linewidth}X@{}}
\toprule
資料 & 位置づけ & 種別 & 時系列 \\
\midrule
MiniMax M3 & 主要資料 & モデル & 2026-06-01 (モデル\_RELEASE) \\
Anthropic June 2026 technical releases & 主要資料 & モデル更新 & 2026-06-30 (モデル\_RELEASE); 2026-06-30 (PRODUCT\_RELEASE); 2026-06-30 (SAFETY\_UPDATE) \\
Alibaba Model Studio June 2026 model lifecycle & 主要資料 & モデル更新 & 2026-06-01 (モデル\_RELEASE); 2026-06-10 (モデル\_UPDATE); 2026-06-22 (モデル\_RELEASE); 2026-06-25 (モデル\_RELEASE); 2026-06-26 (モデル\_AVAILABILITY) \\
Gemini API June 2026 lifecycle & 主要資料 & モデル更新 & 2026-06-30 (モデル\_PREVIEW); 2026-06-30 (モデル\_RELEASE); 2026-06-24 (TOOL\_PREVIEW); 2026-06-17 (API\_UPDATE); 2026-06-15 (DEPRECATION); 2026-06-01 (モデル\_SHUTDOWN) \\
Kimi Code June 2026 releases & 主要資料 & モデル更新 & 2026-06-02 (AGENT\_FEATURE); 2026-06-09 (MULTI\_AGENT\_FEATURE); 2026-06-12 (モデル\_RELEASE) \\
GPT-5.6 series limited preview & 補足資料 & モデル & 2026-06-26 (LIMITED\_PREVIEW) \\
Google DeepMind June 2026 model and agent-safety releases & 補足資料 & モデル更新 & 2026-06-10 (モデル\_RELEASE); 2026-06-18 (SAFETY\_研究); 2026-06 (モデル\_RELEASE) \\
\bottomrule
\end{tabularx}
\end{center}
\normalsize

\begin{technicalnote}{MiniMax M3}{主要資料}
\begin{tabularx}{\linewidth}{@{}>{\bfseries}p{0.22\linewidth}X@{}}
組織 & MiniMax \\
種別 & モデル \\
時系列 & 2026-06-01 (モデル\_RELEASE) \\
\end{tabularx}
\smallskip
{\bfseries 一次資料から整理したtechnical points}
\begin{itemize}[leftmargin=1.5em,itemsep=0.35em]
\item \textbf{Vendor claim}: MiniMax officially released M3 on June 1, describing MSA sparse attention, up to 1M context, native image/video input and desktop-computer operation, with open weights.
\end{itemize}
{\bfseries 読む際の境界}
\begin{itemize}[leftmargin=1.5em,itemsep=0.35em]
\item \textbf{分析上の留意点}: This card is derived from the collected first-party index/release-note snapshot. Capability or benchmark language on vendor pages is not independently reproduced.
\end{itemize}
{\bfseries 一次資料}
\begingroup\sloppy
\begin{itemize}[leftmargin=1.5em,itemsep=0.25em]
\item {\scriptsize\url{https://www.minimax.io/news}}
\end{itemize}
\endgroup
\end{technicalnote}
\medskip

\begin{technicalnote}{Anthropic June 2026 technical releases}{主要資料}
\begin{tabularx}{\linewidth}{@{}>{\bfseries}p{0.22\linewidth}X@{}}
組織 & Anthropic \\
種別 & モデル更新 \\
時系列 & 2026-06-30 (モデル\_RELEASE); 2026-06-30 (PRODUCT\_RELEASE); 2026-06-30 (SAFETY\_UPDATE) \\
\end{tabularx}
\smallskip
{\bfseries 一次資料から整理したtechnical points}
\begin{itemize}[leftmargin=1.5em,itemsep=0.35em]
\item \textbf{Vendor claim}: Anthropic's June 30 newsroom records Claude Sonnet 5, Claude Science, and the Fable 5 redeployment/safety-framework announcement.
\end{itemize}
{\bfseries 読む際の境界}
\begin{itemize}[leftmargin=1.5em,itemsep=0.35em]
\item \textbf{分析上の留意点}: This card is derived from the collected first-party index/release-note snapshot. Capability or benchmark language on vendor pages is not independently reproduced.
\end{itemize}
{\bfseries 一次資料}
\begingroup\sloppy
\begin{itemize}[leftmargin=1.5em,itemsep=0.25em]
\item {\scriptsize\url{https://www.anthropic.com/news}}
\end{itemize}
\endgroup
\end{technicalnote}
\medskip

\begin{technicalnote}{Alibaba Model Studio June 2026 model lifecycle}{主要資料}
\begin{tabularx}{\linewidth}{@{}>{\bfseries}p{0.22\linewidth}X@{}}
組織 & Alibaba Cloud \\
種別 & モデル更新 \\
時系列 & 2026-06-01 (モデル\_RELEASE); 2026-06-10 (モデル\_UPDATE); 2026-06-22 (モデル\_RELEASE); 2026-06-25 (モデル\_RELEASE); 2026-06-26 (モデル\_AVAILABILITY) \\
\end{tabularx}
\smallskip
{\bfseries 一次資料から整理したtechnical points}
\begin{itemize}[leftmargin=1.5em,itemsep=0.35em]
\item \textbf{Vendor claim}: Alibaba Cloud's June lifecycle log records Qwen3.7-Plus, a Qwen3.7-Max snapshot with added visual understanding, HappyHorse1.1 video generation, Qwen Image 2.0 Pro, and Singapore-region availability for Kimi K2.7 Code.
\end{itemize}
{\bfseries 読む際の境界}
\begin{itemize}[leftmargin=1.5em,itemsep=0.35em]
\item \textbf{分析上の留意点}: This card is derived from the collected first-party index/release-note snapshot. Capability or benchmark language on vendor pages is not independently reproduced.
\end{itemize}
{\bfseries 一次資料}
\begingroup\sloppy
\begin{itemize}[leftmargin=1.5em,itemsep=0.25em]
\item {\scriptsize\url{https://www.alibabacloud.com/help/en/model-studio/newly-released-models}}
\end{itemize}
\endgroup
\end{technicalnote}
\medskip

\begin{technicalnote}{Gemini API June 2026 lifecycle}{主要資料}
\begin{tabularx}{\linewidth}{@{}>{\bfseries}p{0.22\linewidth}X@{}}
組織 & Google \\
種別 & モデル更新 \\
時系列 & 2026-06-30 (モデル\_PREVIEW); 2026-06-30 (モデル\_RELEASE); 2026-06-24 (TOOL\_PREVIEW); 2026-06-17 (API\_UPDATE); 2026-06-15 (DEPRECATION); 2026-06-01 (モデル\_SHUTDOWN) \\
\end{tabularx}
\smallskip
{\bfseries 一次資料から整理したtechnical points}
\begin{itemize}[leftmargin=1.5em,itemsep=0.35em]
\item \textbf{Vendor claim}: The Gemini API June changelog adds Gemini Omni Flash preview and Gemini 3.1 Flash Lite Image GA on June 30, Computer Use for Gemini 3.5 Flash on June 24, streaming speech generation on June 17, and several lifecycle shutdown/deprecation events.
\end{itemize}
{\bfseries 読む際の境界}
\begin{itemize}[leftmargin=1.5em,itemsep=0.35em]
\item \textbf{分析上の留意点}: This card is derived from the collected first-party index/release-note snapshot. Capability or benchmark language on vendor pages is not independently reproduced.
\end{itemize}
{\bfseries 一次資料}
\begingroup\sloppy
\begin{itemize}[leftmargin=1.5em,itemsep=0.25em]
\item {\scriptsize\url{https://ai.google.dev/gemini-api/docs/changelog}}
\end{itemize}
\endgroup
\end{technicalnote}
\medskip

\begin{technicalnote}{Kimi Code June 2026 releases}{主要資料}
\begin{tabularx}{\linewidth}{@{}>{\bfseries}p{0.22\linewidth}X@{}}
組織 & Moonshot AI / Kimi \\
種別 & モデル更新 \\
時系列 & 2026-06-02 (AGENT\_FEATURE); 2026-06-09 (MULTI\_AGENT\_FEATURE); 2026-06-12 (モデル\_RELEASE) \\
\end{tabularx}
\smallskip
{\bfseries 一次資料から整理したtechnical points}
\begin{itemize}[leftmargin=1.5em,itemsep=0.35em]
\item \textbf{Vendor claim}: Kimi Code's changelog records autonomous Goal mode on June 2, /swarm multi-agent parallel tasks on June 9, and the released/open-sourced Kimi K2.7 Code model on June 12.
\end{itemize}
{\bfseries 読む際の境界}
\begin{itemize}[leftmargin=1.5em,itemsep=0.35em]
\item \textbf{分析上の留意点}: This card is derived from the collected first-party index/release-note snapshot. Capability or benchmark language on vendor pages is not independently reproduced.
\end{itemize}
{\bfseries 一次資料}
\begingroup\sloppy
\begin{itemize}[leftmargin=1.5em,itemsep=0.25em]
\item {\scriptsize\url{https://www.kimi.com/code/docs/en/kimi-code/whats-new.html}}
\end{itemize}
\endgroup
\end{technicalnote}
\medskip

\begin{technicalnote}{GPT-5.6 series limited preview}{補足資料}
\begin{tabularx}{\linewidth}{@{}>{\bfseries}p{0.22\linewidth}X@{}}
組織 & OpenAI \\
種別 & モデル \\
時系列 & 2026-06-26 (LIMITED\_PREVIEW) \\
\end{tabularx}
\smallskip
{\bfseries 一次資料から整理したtechnical points}
\begin{itemize}[leftmargin=1.5em,itemsep=0.35em]
\item \textbf{Vendor claim}: OpenAI began a limited preview of GPT-5.6 Sol, Terra and Luna on June 26; Sol was limited to a small group of trusted partners, with broader availability planned later.
\end{itemize}
{\bfseries 読む際の境界}
\begin{itemize}[leftmargin=1.5em,itemsep=0.35em]
\item \textbf{分析上の留意点}: The source is first-party OpenAI material. Capability, benchmark, efficiency, or performance statements remain vendor claims unless explicitly described as externally reproduced.
\end{itemize}
{\bfseries 一次資料}
\begingroup\sloppy
\begin{itemize}[leftmargin=1.5em,itemsep=0.25em]
\item {\scriptsize\url{https://openai.com/index/previewing-gpt-5-6-sol}}
\end{itemize}
\endgroup
\end{technicalnote}
\medskip

\begin{technicalnote}{Google DeepMind June 2026 model and agent-safety releases}{補足資料}
\begin{tabularx}{\linewidth}{@{}>{\bfseries}p{0.22\linewidth}X@{}}
組織 & Google DeepMind \\
種別 & モデル更新 \\
時系列 & 2026-06-10 (モデル\_RELEASE); 2026-06-18 (SAFETY\_研究); 2026-06 (モデル\_RELEASE) \\
\end{tabularx}
\smallskip
{\bfseries 一次資料から整理したtechnical points}
\begin{itemize}[leftmargin=1.5em,itemsep=0.35em]
\item \textbf{Vendor claim}: Google DeepMind's June index includes DiffusionGemma, Gemma 4 12B, computer-use/voice model updates, and agent-safety research.
\end{itemize}
{\bfseries 読む際の境界}
\begin{itemize}[leftmargin=1.5em,itemsep=0.35em]
\item \textbf{分析上の留意点}: This card is derived from the collected first-party index/release-note snapshot. Capability or benchmark language on vendor pages is not independently reproduced.
\end{itemize}
{\bfseries 一次資料}
\begingroup\sloppy
\begin{itemize}[leftmargin=1.5em,itemsep=0.25em]
\item {\scriptsize\url{https://deepmind.google/blog/}}
\end{itemize}
\endgroup
\end{technicalnote}
\medskip
